\begin{center}
\textbf{\large{Abstract}}
\end{center}

\addcontentsline{toc}{chapter}{Abstract}
\begin{small}

Access control restricts the operations that each user is allowed to perform on a REST API. Current automated testing techniques mainly target stateless authorization flaws, such as object-level (BOLA) or function-level (BFLA) flaws. Stateful broken access control is rarely checked directly by these techniques. This flaw occurs when a user's privileges are revoked at runtime, but the system keeps accepting the user's existing sessions.

This thesis proposes an automated approach to detect this class of vulnerabilities through metamorphic testing in the EvoMaster tool. The approach is based on the metamorphic relation MR-A1: an operation that succeeded before the account was downgraded must be rejected when it is sent again later with the same token. Each test scenario consists of an operation triple \texttt{[source/transition/follow-up]}. A Naive Bayes model automatically identifies downgrade operations from the OpenAPI specification, which reduces manual configuration.

The experiments comprise 98 black-box runs on four target applications. Two of them, reports-api and vampi-seeded, carry seeded bugs. pyLoad carries the vulnerability CVE-2026-41133, and artifact-keeper carries a bug that its developers have since fixed. On the vulnerable variants, the tool detected violations in 19/20 runs on reports-api, 7/20 runs on vampi-seeded, 5/5 runs on pyLoad and 5/5 runs on artifact-keeper. Across 48 runs on the fixed variants, including the developers' own fix of artifact-keeper, the tool raised no false alarms. MR-A1 detected servers that keep honouring revoked privileges in all four studied applications.

\vspace*{1cm}
\textbf{Keywords}: REST API, access control, stateful authorization, metamorphic testing, EvoMaster
\end{small}
